%! Equivalent Representations of Polynomial Expressions — Unit 2, Lesson 2.6 — Group Activity
\documentclass[10pt]{article}
\usepackage{precalculus-article}
\usepackage{precalculus-boxes}
\newcommand{\TallMath}[1]{$\displaystyle #1\rule[-1.4em]{0pt}{3.2em}$}

\begin{document}

\pageheader{Unit 2, Lesson 2.6}{Group Activity}

\vspace{0.12in}

\begin{scenariobox}[Same Polynomial, Two Jobs]{plumlight}
\small
Your group is the conversion desk. Every polynomial that arrives has a
question attached, and the job is to get it into the form that answers the
question --- and no further. Two references, good for every tier:

\smallskip
\begin{center}
\renewcommand{\arraystretch}{1.1}
\begin{tabular}{r c}
row $0$: & $1$ \\
row $1$: & $1 \quad 1$ \\
row $2$: & $1 \quad 2 \quad 1$ \\
row $3$: & $1 \quad 3 \quad 3 \quad 1$ \\
row $4$: & $1 \quad 4 \quad 6 \quad 4 \quad 1$ \\
row $5$: & $1 \quad 5 \quad 10 \quad 10 \quad 5 \quad 1$ \\
\end{tabular}
\end{center}

\smallskip
\textbf{Division algorithm:} $f(x) = g(x)\,q(x) + r(x)$ with
$\deg r < \deg g$.
\end{scenariobox}

\vspace{0.1in}

\boxguard
\begin{tcolorbox}[colback=white, colframe=black!40,
    title=\textbf{Tier R --- Remediate},
    fonttitle=\bfseries, arc=2mm, left=3mm, right=3mm, top=2mm, bottom=2mm]

\begin{enumerate}[itemsep=8pt]
  \item Write \textbf{factored} or \textbf{standard} --- whichever form
        answers the question faster.

        \begin{enumerate}[label=(\alph*), itemsep=4pt]
          \item Where does the graph cross the $x$-axis? \quad \blank{2.6cm}
          \item What is $\lim_{x \to \infty} p(x)$? \quad \blank{2.6cm}
          \item What is the $y$-intercept? \quad \blank{2.6cm}
          \item What are the real zeros? \quad \blank{2.6cm}
        \end{enumerate}

  \item Row 3 of Pascal's Triangle is \blank{1cm}, \blank{1cm}, \blank{1cm},
        \blank{1cm}. Use it to expand $(x+1)^3$.

        \begin{work}
          (x+1)^3 &= x^3 + 3x^2(1) + 3x(1)^2 + (1)^3 \\
                  &= x^3 + 3x^2 + 3x + 1
        \end{work}

        $(x+1)^3 = $ \blank{6cm}

  \item Multiply out $(x-5)(x+2)$ and write the result in standard form.

        \begin{work}
          (x-5)(x+2) &= x^2 + 2x - 5x - 10 \\
                     &= x^2 - 3x - 10
        \end{work}

        $(x-5)(x+2) = $ \blank{5cm}

  \item Complete the division algorithm.

        \begin{enumerate}[label=(\alph*), itemsep=4pt]
          \item $f(x) = g(x)\,\cdot$ \blank{1.6cm} $+$ \blank{1.6cm}
          \item The degree of $r$ must be \blank{2.4cm} the degree of $g$.
          \item If $g$ is a factor of $f$, the remainder is \blank{1.4cm}.
        \end{enumerate}
\end{enumerate}
\end{tcolorbox}

\vspace{0.1in}

\boxguard
\begin{tcolorbox}[colback=white, colframe=black!40,
    title=\textbf{Tier A --- Approaching Proficiency},
    fonttitle=\bfseries, arc=2mm, left=3mm, right=3mm, top=2mm, bottom=2mm]

\begin{enumerate}[itemsep=8pt]
  \item Let $p(x) = (x+2)(x-1)(x-4)$.

        \begin{enumerate}[label=(\alph*), itemsep=4pt]
          \item Real zeros, straight from this form: \blank{4.4cm}
          \item Write $p$ in standard form.

                \begin{work}
                  p(x) &= (x+2)(x-1)(x-4) \\
                       &= \big(x^2 + x - 2\big)(x-4) \\
                       &= x^3 - 4x^2 + x^2 - 4x - 2x + 8 \\
                       &= x^3 - 3x^2 - 6x + 8
                \end{work}

                $p(x) = $ \blank{6cm}
          \item $y$-intercept: \blank{1.6cm} \quad Which form gave it to you?
                \blank{2.4cm}
          \item $\lim_{x \to -\infty} p(x) = $ \blank{1.6cm} \quad Which form?
                \blank{2.4cm}
        \end{enumerate}

  \item Use row 4 of Pascal's Triangle to expand $(x+3)^4$. Do not multiply
        four binomials.

        \begin{work}
          (x+3)^4 &= x^4 + 4x^3(3) + 6x^2(3)^2 + 4x(3)^3 + (3)^4 \\
                  &= x^4 + 12x^3 + 54x^2 + 108x + 81
        \end{work}

        $(x+3)^4 = $ \blank{7cm}

  \item Divide $x^3 + 2x^2 - 5x - 6$ by $x + 1$.

        \begin{work}
          &\begin{array}{r@{\;}rrrr}
                      &   x^2 &   +x  &  -6  &     \\ \cline{2-5}
           x+1\,\big) &   x^3 & +2x^2 & -5x  & -6  \\
           -          &   x^3 &  +x^2 &      &     \\ \cline{2-5}
                      &       &   x^2 & -5x  & -6  \\
           -          &       &   x^2 &  +x  &     \\ \cline{2-5}
                      &       &       & -6x  & -6  \\
           -          &       &       & -6x  & -6  \\ \cline{2-5}
                      &       &       &      &  0
          \end{array}
        \end{work}

        Quotient: \blank{3.6cm} \qquad Remainder: \blank{1.2cm}

        Is $x+1$ a factor of $x^3 + 2x^2 - 5x - 6$? \quad \blank{1.6cm}
\end{enumerate}
\end{tcolorbox}

\vspace{0.1in}

\boxguard
\begin{tcolorbox}[colback=white, colframe=black!40,
    title=\textbf{Tier E --- Extension},
    fonttitle=\bfseries, arc=2mm, left=3mm, right=3mm, top=2mm, bottom=2mm]

\begin{enumerate}[itemsep=8pt]
  \item Someone claims $(x-2)$ is a factor of $f(x) = x^3 - 6x^2 + 11x - 6$.
        Test the claim by dividing.

        \begin{work}
          &\begin{array}{r@{\;}rrrr}
                      &   x^2 &  -4x  &  +3  &     \\ \cline{2-5}
           x-2\,\big) &   x^3 & -6x^2 & +11x & -6  \\
           -          &   x^3 & -2x^2 &      &     \\ \cline{2-5}
                      &       & -4x^2 & +11x & -6  \\
           -          &       & -4x^2 &  +8x &     \\ \cline{2-5}
                      &       &       &  +3x & -6  \\
           -          &       &       &  +3x & -6  \\ \cline{2-5}
                      &       &       &      &  0
          \end{array}
        \end{work}

        \begin{enumerate}[label=(\alph*), itemsep=4pt]
          \item Remainder: \blank{1.2cm} \quad So the claim is
                \textbf{true} / \textbf{false}.
          \item Factor the quotient, then write $f$ completely factored:
                \blank{6cm}
          \item All real zeros of $f$: \blank{3.6cm} \quad
                $y$-intercept: \blank{1.4cm}
        \end{enumerate}

  \item Divide $x^3 + 3x^2 - 2$ by $x^2 - 1$. The divisor is now
        \textit{quadratic} --- and the dividend is missing its $x$ term.

        \begin{work}
          &\begin{array}{r@{\;}rrrr}
                        &       &       &    x  &  +3  \\ \cline{2-5}
           x^2-1\,\big) &   x^3 & +3x^2 &  +0x  &  -2  \\
           -            &   x^3 &       &   -x  &      \\ \cline{2-5}
                        &       & +3x^2 &   +x  &  -2  \\
           -            &       & +3x^2 &       &  -3  \\ \cline{2-5}
                        &       &       &   +x  &  +1
          \end{array}
        \end{work}

        \begin{enumerate}[label=(\alph*), itemsep=4pt]
          \item $q(x) = $ \blank{2.6cm} \qquad $r(x) = $ \blank{2.6cm}
          \item Write the division-algorithm sentence:

                \smallskip
                $x^3 + 3x^2 - 2 = (x^2-1)\,\cdot$ \blank{2.4cm} $+$
                \blank{2.4cm}
          \item $\deg r = $ \blank{1cm} and $\deg g = $ \blank{1cm}, so the
                degree condition \textbf{is} / \textbf{is not} satisfied.
        \end{enumerate}

  \item Explain why ``the remainder is $0$'' and ``the divisor is a factor''
        say exactly the same thing. Use the division algorithm in your answer,
        and connect it to the factor/zero equivalence from Lesson 2.3.

        \vspace{0.05in}
        \writelines{3}
\end{enumerate}
\end{tcolorbox}

\end{document}
